\section{Geometry of the voltage cover and torsion-freeness}
\label{geo:section}

Embedding the components of the voltage cover in a Cayley graph is not
enough for the later operator comparison, because different copies can
overlap.  We therefore seek disjoint assigned domains, with the removed
vertices controlled by one finite ball per copy and the unused Cayley
edges forming a forest.  The input is the girth condition
\eqref{geo:girth}, which the later sampling argument must establish;
no assertion about an operator's spectral distribution is used here.
The constants in this section are chosen before the voltage values and
the final alphabet are chosen.

The use of labeled graphs as relators follows the geometric strategy of
Gromov \cite[Sections~2 and~4.8]{Gromov2003}. Finite alphabets with global
separation of long path words are developed by Osajda
\cite[Section~2]{Osajda2020}.  The diagrammatic embedding argument is
closely related to Ollivier \cite[Section~2]{Ollivier2006} and Gruber
\cite[Lemma~4.1]{Gruber2015}.
Here the relator graphs are components of a voltage cover, possibly with
nontrivial stabilizers. We give the embedding, overlap, and torsion
arguments for these precise hypotheses.

\subsection{Graphs, presentations, and the conditional input}
\label{geo:conventions}

All graphs in this section are combinatorial graphs, possibly with loops or
parallel edges.  An unoriented edge has two formal orientations, interchanged
by a fixed-point-free reversal operation.  Traversing an edge and immediately
traversing its reverse is a backtrack.  A closed walk is cyclically
nonbacktracking if it has no backtrack, including at its closing corner.
Edge lengths are one, degrees count incidences, and a ball denotes the induced
subgraph on the vertices at most its stated distance from its center.  A
forest has no nonempty finite circuit, with a loop and a pair of parallel
edges counted as circuits of lengths one and two.  Chains always have finite
support.  An orientation chosen once for each unoriented edge identifies the
space of edge chains over a field $k$ with the free vector space on those
edges, with boundary $\partial[x,y]=[y]-[x]$.  The cycle space is
$Z_1(-;k)=\ker\partial$.

Let $Y$ be the surviving tower.  The properties of it needed in this section
are the following: it is countable, has maximum degree at most $6D$, and has
an integer layer function whose change across an edge is at most one.  Its
oriented labeling is folded: the labels of the oriented edges leaving any
vertex are distinct.  Positive labels and their formal inverses are disjoint
alphabets.  Finally a nonbacktracking path starting in layer $i$, of length
at least $L_i=100H_i/\log_2d$, is determined, including its location in $Y$,
by its oriented label word.  This uniqueness compares paths anywhere in $Y$.
Refining these labels by arbitrary additional tags preserves both foldedness
and this uniqueness property.

The following elementary estimate will be used with fixed constants, however
large.  If two vertices can be joined by a walk of length
$aH_i/\log_2d$, starting in layer $i$, then
\begin{equation}
 \left|\frac{H_j}{H_i}-1\right|\le\frac{a}{\log_2d}
 \label{geo:layer-comparison}
\end{equation}
at every layer $j$ encountered.  This follows from
$|H_j-H_i|\le |j-i|$.  Consequently all subsequent comparisons become uniform
in $i$ by increasing $\ell$.  Also $L_i>2$ for all $i$ with the parameters
of the construction.

Let $F$ be a countable torsion-free group.  Give each positive edge $e$ of
$Y$ a voltage $z_e\in F$, and give its reverse voltage $z_e^{-1}$.  The ordinary voltage-cover construction of Gross and Gross--Tucker
\cite{Gross1974,GT1977} gives a graph $P$ with vertices $(v,f)$ and an oriented lift
\[
 (v,f)\longrightarrow(v',fz_e)
\]
over each oriented edge $e:v\to v'$.  The left deck action is
$\delta_a(v,f)=(v,af)$ for $a\in F$.  It acts freely on vertices and
transitively on the lifts of every fixed base vertex or fixed base path;
it can carry one connected component of $P$ to another.  Labels, layers,
and the distinguished positive orientations are pulled back from $Y$.

We assume that the augmented positive alphabet $S$ is finite and that
\begin{equation}
 \begin{gathered}
 \text{there is no nonempty cyclically nonbacktracking closed walk in $P$}\\
 \text{based in layer $i$ of length at most }
             C_*H_i/\log_2d,\qquad i\in\mathbb Z.
 \end{gathered}
 \label{geo:girth}
\end{equation}
Here $C_*$ is a single absolute constant fixed below.  This is a condition
on the actual cover, not just on $Y$ and not just on simple base circuits.
It includes all lifts, although whether a base closed walk has trivial
holonomy does not depend on the starting deck coordinate.

Define
\begin{equation}
 \Gbase=\langle S\mid \text{label words of all closed walks in }P\rangle,
 \qquad X=\operatorname{Cay}_{\mathrm{right}}(\Gbase,S).
 \label{geo:presentation}
\end{equation}
The Cayley graph is formal: its positive edges are the pairs $(h,\sigma)$,
with endpoints $h,h\sigma$, for $h\in\Gbase$, $\sigma\in S$.  We do not
identify distinct such edges even if their unordered endpoints coincide.
Their reverse orientations have labels $\sigma^{-1}$.  Left translation
preserves every formal edge label and its positive orientation.

For a connected component $C$ of $P$, a vertex $v_0\in C$, and $h\in\Gbase$,
there is a unique label-preserving graph map $\phi:C\to X$ with
$\phi(v_0)=h$: the image of the endpoint of a path from $v_0$ is $h$ times
the label word of that path.  Independence of the path follows from
\eqref{geo:presentation}.  The maps are not assumed injective at this stage.

\begin{theorem}[Geometric conclusion]
\label{geo:main}
There are absolute constants $C_*,C_\bullet$ such that, for sufficiently
large $\ell$, the preceding hypotheses have the following consequences.
Every component map $\phi:C\to X$ is an embedding.  Index copies by pairs
$(C,\phi)$ modulo the equivalence
\[
 (C,\phi)\sim
 (\delta_a C,\phi\delta_a^{-1}),\qquad a\in F.
\]
Thus an indexed copy $Q$ remembers its source component and source vertices
up to this deck alignment, in addition to its embedded image.
There is a well-order of these copies and an induced intrinsic ball $B(Q)$
in each $Q$ such that
\begin{enumerate}
 \item every vertex of $Q$ shared with an earlier copy belongs to $B(Q)$;
 \item the subgraph of $Q$ consisting of all edges with at least one
 endpoint in $B(Q)$ is a forest;
 \item for every vertex of $B(Q)$ lying over layer $i$,
 \begin{equation}
 |V(B(Q))|\le\mathfrak w_i,
 \qquad \mathfrak w_i=2^{C_\bullet H_i};
 \label{geo:weight}
 \end{equation}
 \item if $D_Q$ is $V(Q)$ minus its vertices shared with earlier copies,
 the domains $D_Q$ are pairwise disjoint, and deleting every edge of $Q$
 whose two endpoints lie in $D_Q$, for all $Q$, leaves a forest in $X$.
\end{enumerate}
Furthermore $\Gbase$ is torsion-free.  None of the constants depends on
$F$, the voltage values, the number of components of $P$, or $|S|$.
\end{theorem}

\subsection{Plane diagrams and reductions before injectivity}
\label{geo:diagrams}

The same short-arc count will first rule out failures of injectivity and
later bound intrinsic distances within embedded copies.  We retain
source paths in the diagrams: merging source walks requires actual deck
alignment, not merely equality of label words.

We use the following form of the van Kampen Lemma; see
\cite[Chapter V, Section 1]{LyndonSchupp1977}.  If a word is freely equal to a
product of $J$ conjugates of defining relators and their inverses, it is
the exterior contour word of a finite connected plane labeled multigraph
with at most $J$ bounded faces.  Each bounded face contour reads a defining
relator, with a chosen cyclic start and orientation.  Conversely an exterior
contour word of such a diagram is freely a product of conjugates of its
bounded face words.  In this statement a contour is a closed walk around a
face, rather than an embedded circle: vertices and bridges may occur more
than once.  Every application below uses a finite number of relators even
though the presentation is infinite.

For completeness, the elementary operations in this version are as follows.
Draw a polygon for each relation factor and its conjugating stem, joining
the stems at a base vertex in their product order.  Empty factors are
omitted.  The exterior reads the literal product.  Exterior insertion of
$aa^{-1}$ is made by attaching a spur.  For a cancellation, let the two
consecutive exterior traversals have vertices $v,w,v'$.  If they traverse
the same edge, erase the resulting leaf.  If they are distinct and $v=v'$,
their digon, or the two lobes when both edges are loops, bounds portions
which can be discarded, together with these edges.  The other exterior
contour remains connected at $v$.  The consecutiveness at $w$ ensures that
no exterior attachment is lost away from that vertex.  If $v\ne v'$, draw
an auxiliary edge in the exterior face between the indicated occurrences
of $v$ and $v'$, so that it cuts off the two-edge portion as a triangular
face.  Contract the auxiliary edge.  The artificial face is now a digon;
delete one of its edges to merge it with the adjacent face.  In that
adjacent contour the replacement traversal has the same label.  Thus the
desired exterior cancellation has taken place without increasing the
number of relation faces.  This operation uses positions on contours and
is valid when other vertices or contour portions repeat.  Conversely,
merging an exterior-adjacent bounded face across a non-bridge edge,
repeatedly until a tree remains, proves the product assertion.  These
operations allow a prescribed literal exterior word to be constructed,
after which its required break positions are marked.  A marked position
crossed by an intermediate cancellation need not retain its previous
vertex value.  All vertices of a diagram have consistent values in $\Gbase$:
propagate from one vertex along paths; simple circuits bound subdiagrams,
so the product assertion proves path independence.

In our presentation every face comes with an actual closed walk in $P$
reading its contour.  This additional assignment matters: equal words alone
will not be treated as equal paths.  When two source incidences have been
aligned by a deck transformation, deleting an edge between their faces
splices two genuine source walks into a genuine closed source walk.  This
is again a permitted relator because the presentation uses all closed
walks, including walks with backtracks and singular contours.

Suppose an injectivity failure exists.  Choose, over \emph{all} component
maps, pairs of distinct source vertices with equal image, paths between
them, and diagrams for the resulting closed image word, a choice minimizing
first the number $J$ of bounded faces and then the number of edges.  The
exterior has one adjustable side: it is this source path, with its break
vertex marked.  Its endpoints in $P$ must remain distinct, but may vary
when comparing different failure diagrams.  Minimization over all failures,
rather than just one originally selected path, is essential below.

We will use the same minimization after injectivity is proved, in a second
setting.  Adjoin to $X$ one center $c_Q$ for every indexed copy $Q$, and one
spoke of length one from $c_Q$ to each vertex of $Q$.  Denote this graph
by $\widehat X$.  For an edge loop in $\widehat X$, started at a group
vertex, each generator edge is a fixed side and each consecutive two-spoke
passage through a center is an adjustable side.  Inflate an adjustable
side to any source path in its copy with the indicated endpoints.  Mark
all side breaks.  If there are $m$ sides there are at most $m$ distinct
marked vertices.  Among diagrams and these adjustable paths minimize
$(J,\text{number of edges})$ lexicographically.  If the inflation is empty,
retain its one marked vertex.  For both minimizations, source data along
adjustable sides and bounded contours are retained as part of the diagram.

\begin{lemma}[Reduction rules]
\label{geo:reductions}
In either minimum diagram the following statements hold.
\begin{enumerate}
 \item No bounded face has a freely trivial label word.
 \item No edge separating two distinct bounded faces has aligned source
 incidences.  Nor does an edge between a bounded face and an adjustable
 exterior side have aligned source incidences.
 \item There is no unmarked vertex of degree one, and no bridge having
 both incidences on a bounded face.
\end{enumerate}
Here alignment means that, when both incidences are read in the same
direction, a left deck transformation identifies their source edges.
\end{lemma}

\begin{proof}
A freely trivial face word can be omitted from the face product expression;
the van Kampen Lemma then supplies a diagram for the identical exterior
with fewer bounded faces.  For the second assertion, delete the separating
edge.  Distinct bounded faces merge into a face whose source walk is obtained
by splicing after deck alignment.  A bounded face adjacent to an adjustable
side can instead be absorbed into that side: replace its edge traversal by
the complementary face walk with the appropriate orientation.  The source
endpoints are unchanged.  Deleting an edge separating two faces does not
disconnect the plane graph, and repeated contacts elsewhere merely make
the new contour singular; they do not prevent the splice.  In each case
$J$ decreases.

At an unmarked leaf the contour makes a backtrack.  If it is an exterior
corner it lies within an adjustable side, since a break between fixed sides
would be marked.  In all cases foldedness makes it an actual backtrack in
the assigned source walk.  Erasing the leaf preserves the assignments and
contradicts edge minimality.

Consider now a bridge with both incidences on a bounded face.  Removing its
open edge splits the diagram in two.  The exterior contour and all its
marked positions belong to one of them; retain that component.  In the
affected bounded face contour the other component is visited in one
excursion, along the bridge, around its exterior contour, and back.  The
discarded component, with the bridge as a stem, gives a diagram for the
word of this excursion with fewer than $J$ bounded faces: the affected
face is not among them.  To see the indicated exterior designation, the
retained component is connected and accessible from infinity, so lies on
the unbounded-contour side of the discarded component.  Thus its boundary
is exactly the contour traversed in that excursion, with the bridge stem.

After injectivity, the two endpoints of the source excursion are equal
because their images are equal, so the excursion can be removed.  Before
injectivity, if those source endpoints were distinct, this excursion and
its smaller diagram would themselves constitute an injectivity failure,
contrary to the global minimal choice of $J$.  They therefore agree in
this case as well.  Delete the excursion, bridge, and discarded component.
The retained affected face is still assigned a closed source walk.  Its
face count does not increase and its edge count decreases, a contradiction.
\end{proof}

Suppress every unmarked degree-two vertex.  Call a resulting edge an
\emph{arc}; it remembers its entire original edge path and its length.
Loops among these arcs are allowed.  There is at least one marked vertex,
so an entire degree-two component cannot disappear without a retained
vertex.  Every side incidence along an arc is constant: it is a bounded
face or one exterior side.  Indeed a change of exterior side is marked,
and a change of face incidence requires branching.  Degrees and contours
in what follows are those of this compressed plane graph.

\begin{lemma}[Short arcs and Euler count]
\label{geo:arcs}
Every arc incident with a bounded face has length less than $L_i$, when
read from an endpoint in layer $i$ on that face.  Its length is also less
than $L_i$ on any adjustable-side incidence with which it is compared.
Every bounded face has at least twenty arc incidences.  If $b$ is the
number of arcs and there are at most $m$ marked vertices, then
\begin{equation}
 b\le 3J+3m-3,\qquad 20J\le 2b,\qquad b\le 6m.
 \label{geo:euler}
\end{equation}
These statements require only $C_*\ge10^4$ and sufficiently large $\ell$.
\end{lemma}

\begin{proof}
An arc of length greater than two is freely reduced.  Otherwise an inverse
pair with an unmarked bivalent midpoint can be erased, identifying its two
outer vertices.  Those outer vertices are distinct in the diagram: the
interior of a maximal arc has no repetitions, and a loop arc of length
greater than two has this property for each such pair.  The two-edge path
is therefore a tree, and its contraction preserves the number of faces
and the plane embedding.  Its outer group values agree; its bivalent
middle vertex has no other incidences.  On every assigned source walk
foldedness identifies the inverse pair as a backtrack.  The contraction
therefore preserves all assignments and decreases the number of edges,
which is impossible.

If an arc has a generator-side incidence it has length one.  Otherwise an
arc incident with a bounded face has two source assignments.  Its other
incidence cannot be on the same bounded face, since an edge with the same
face on both sides is a bridge, excluded by Lemma~\ref{geo:reductions}.
If the arc were at least $L_i$ long on either compared assignment, its
freely reduced word would determine its projected nonbacktracking path in
$Y$.  Both projections would be this same path.  The two lifts of that
path then differ by a single deck transformation, so their source
incidences are aligned.  Lemma~\ref{geo:reductions} forbids that alignment.
Arcs of length at most two already satisfy the claimed inequalities.

If a face had fewer than twenty arc incidences, select its longest arc,
of length $a$, and let $i$ be its initial layer in the face assignment.
Its perimeter would be at most $19a<20L_i=2000H_i/\log_2d$.
By \eqref{geo:layer-comparison}, every layer $j$ on it has
$H_j\ge H_i/2$ for large $\ell$.  The face word is not freely trivial.
Cancel backtracks in its closed source walk, and then cancel around its
cyclic ends.  Foldedness makes these actual source cancellations; the
remaining walk is nonempty and cyclically nonbacktracking.  Its length
is at most the displayed perimeter, while
$C_*H_j/\log_2d\ge5000H_i/\log_2d$.  This contradicts
\eqref{geo:girth}.

Let $v$ be the number of vertices of the compressed diagram.  It is
connected and satisfies $b=v+J-1$ by the plane Euler formula, also for
loops, parallel edges and bridges.  Every unmarked vertex has degree at
least three; a graph consisting of one marked vertex and no edges also
satisfies the resulting inequality.  Hence
$2b\ge3(v-m)$, which gives the first inequality of
\eqref{geo:euler}.  Each arc has two face incidences in total, counting
the exterior, so the twenty-incidence bound gives its second inequality.
Combining them gives $b\le(30/7)(m-1)\le6m$ whenever edges occur;
the same final bound holds for the trivial diagram.
\end{proof}

\begin{proposition}[Embeddings]
\label{geo:embeddings}
All component label maps $C\to X$ are injective on vertices and formal
edges.  The indexed copies of Theorem~\ref{geo:main} are therefore defined.
\end{proposition}

\begin{proof}
Apply Lemma~\ref{geo:arcs} to a minimum injectivity-failure diagram, with
$m=1$.  Its first two inequalities imply $20J\le6J$, hence $J=0$.
The uncompressed diagram is a tree, so its exterior word is freely trivial.
A path with freely trivial word in a folded graph returns to its starting
vertex: erase successive inverse pairs, which are actual backtracks.
The chosen source endpoints consequently agree, a contradiction.

For edges, two oriented source edges with the same formal image have the
same initial image and the same oriented label.  Vertex injectivity makes
their initial vertices equal, and foldedness makes the edges equal.
This includes the possibility of parallel edges and establishes the
asserted graph embeddings.
\end{proof}

We henceforth use the post-injectivity minimum diagrams.  One additional
observation follows immediately from the indexing convention.  If paths
in two indexed copies have the same word and are deck-aligned, and their
maps agree at one source point after alignment, their maps agree everywhere
on the connected component: propagate along labeled paths.  They are then
the same indexed copy.  Therefore an arc shared by two different copy
assignments has length less than the applicable $L_i$ on either assignment.
\label{geo:distinct-copy-arcs}

\subsection{Short chain fillings and an explicit thin-triangle estimate}
\label{geo:coned}

We next prepare an ordering of the copy centers by distance from one
fixed group vertex, in order to control all earlier intersections with
a copy at once.  Thin triangles in $\widehat X$ will give uniformly
short paths that avoid the copy's own center.  The short-arc estimate
will then turn such paths into intrinsic distance bounds in the copy.
We first obtain thin triangles from a linear bound on decompositions
of loop chains into short loops.

\begin{lemma}[Short chain fillings]
\label{geo:chain-filling}
In $\widehat X$ the integral oriented chain of any edge loop of length $n$
is a sum of at most $12n$ oriented chains of edge loops of length at most
four.
\end{lemma}

\begin{proof}
For a nonempty loop, rotate its start to a group vertex and split it into
the sides specified above.  There are $m\le n$ sides.  Inflate and minimize,
and apply Lemma~\ref{geo:arcs}.  For each arc choose a standard path between
the images of its endpoints: use the generator edge if it has a fixed-side
incidence, and otherwise use the two-spoke passage through any one of its
incident source-copy centers.  Choose one direction for the arc, using the
reverse standard path for its reverse incidence.

On a bounded face, replace every incident standard path by the two-spoke
passage through that face's center.  Each replacement differs by a closed
chain of length at most four; a generator standard path only reduces that
length.  The sum of the new passages is zero as a chain, because successive
spokes telescope at the consecutive contour vertices, with their contour
multiplicities.  On an exterior adjustable side make the same replacements
using that side's center; the resulting chain telescopes to the original
two-spoke side.  Exterior fixed sides already have the required standard
path.  The sum of the standard-path chains around the bounded faces is
the exterior standard-path chain, by cancellation of the two oriented
incidences of each internal edge, also when an edge is a bridge or a
contour is singular.  Combining these identities expresses the original
loop chain as the asserted sum.  At most one short loop is used for each
bounded or exterior arc incidence, hence at most $2b\le12m\le12n$.
\end{proof}

The next lemma is a quantitative form of the homological-isoperimetric
criterion for hyperbolicity; compare Gersten
\cite[Theorem~6.3]{Gersten1996}.  The proof below gives an explicit
thin-triangle bound under our chain-filling hypothesis, without assuming
local finiteness.

\begin{lemma}[Thin triangles from short chain fillings]
\label{geo:thin}
Let a connected graph with unit edge lengths have the conclusion of
Lemma~\ref{geo:chain-filling}.  Every geodesic triangle on its vertices is
$D_0$-thin, where one may take $D_0=10^6$.  Thinness here measures distance
to the vertices on the other two sides.
\end{lemma}

\begin{proof}
Fix a geodesic triangle $[x,y],[x,w],[y,w]$, and put
\[
 n=d(x,y),\quad u=d(x,w),\quad v=d(y,w),\quad
 a=(n+u-v)/2,\quad b=u-a=v-(n-a).
\]
Let $R$ be the maximum, over its side vertices, of distance to the vertices
on the other two sides.  Relabel the sides so it is attained at
$z\in[x,y]$, at position $j=d(x,z)$.  We may assume $R\ge100$.  Then
$j,n-j\ge R$, and every vertex of either other side has distance at least
$R$ from $z$.

A position $h<a-R$ on $[x,y]$ is within $R$ of $[x,w]$.  Indeed, if it
were within $R$ of position $t$ from $y$ on $[y,w]$, then
$t\ge n-h-R$, and the triangle inequality would give
\[
 u\le h+R+v-t\le2h+2R+v-n,
\]
contrary to $h<a-R$.  The definition of $R$ supplies closeness to at least
one other side.  Symmetrically positions $h>a+R$ are within $R$ of
$[y,w]$.  Also a position $k>a+R$ on $[x,w]$ must be within $R$ of
$[y,w]$: closeness to a position $h$ on $[x,y]$ would imply $h\ge k-R$
and
$v\le u-k+R+n-h\le u+n-2k+2R$, again a contradiction.

We construct a subsegment of $[x,y]$ containing at least $R$ edges on
either side of $z$, and an alternate path between its endpoints whose
vertices all have distance at least $R$ from $z$, so that the resulting
loop has length at most $200R$.

First suppose $j<a-6R$.  Take the positions
$h_- =\max(0,j-4R)$ and $h_+=j+4R$; the latter is less than $a-2R<n$.
Bridge both to $[x,w]$ by paths of length at most $R$, taking the trivial
bridge at $x$ if $h_-=0$, and join their endpoints along $[x,w]$.
Each nontrivial bridge starts $4R$ from $z$, so all its vertices are at
least $3R$ from $z$.  The side portion on $[x,w]$ has length at most
$h_+-h_-+2R$, since its endpoint coordinates from $x$ differ from $h_-$
and $h_+$ by at most $R$.  Thus the whole loop has length at most $20R$.
The case $j>a+6R$ is symmetric.

In the remaining case $|j-a|\le6R$, take
\[
 h_- =\max(0,\lfloor a-10R\rfloor),\qquad
 h_+ =\min(n,\lceil a+10R\rceil).
\]
Bridge the left endpoint to $[x,w]$ and the right endpoint to $[y,w]$,
using paths of length at most $R$ and trivial bridges at clipped endpoints.
The preceding positional rules supply every nontrivial bridge.  Such a
bridge starts at least $4R$ from $z$, so it stays at least $3R$ from $z$.
Write $s_-$ and $t_+$ for the coordinates of the bridge endpoints from
$x$ and $y$, respectively.  They satisfy
\[
 |s_--h_-|\le R,\quad |t_+-(n-h_+)|\le R,\quad
 a-h_-\le10R+1,\quad h_+-a\le10R+1.
\]
If $b\le20R+2$, join these endpoints through $w$.  Each side portion
has length at most $31R+3$, while $h_+-h_-\le20R+2$, so the entire loop
has length at most $84R+8$.  If $b>20R+2$, take the vertex $A\in[x,w]$
at position $k=\lceil a+20R\rceil<u$.  It has a bridge of length at most
$R$ to a vertex $B\in[y,w]$.  If $t=d(y,B)$, comparison of distances to
$w$ gives
\[
 |(u-k)-(v-t)|\le R,\qquad
 |t-(n-a+20R)|\le R+1.
\]
We have $s_-\le k$ and $t_+\le t$; the intervening portions on the two
sides have lengths at most $31R+2$ and $32R+2$.  Including the three
bridges and the original segment gives a loop of length at most $86R+6$.
This last bridge stays far from $z$, since
$d(z,A)\ge k-j\ge14R$ and its length is at most $R$.
All other vertices of the alternate path lie on the two other sides or
the already controlled endpoint bridges.  In all cases the selected
segment extends at least $R$ on either side of $z$: when clipped, this
uses $j,n-j\ge R$; when not clipped it follows from the stated positions.
This completes the construction with the claimed $200R$ bound.

Put $r_0=\lfloor R/4\rfloor$.  In cyclic order partition this loop into
four vertex paths: $I_1$ is the segment from $j-r_0$ to $j+r_0$;
$I_2$ is the remaining segment toward $h_+$; $I_3$ is the alternate path
from $h_+$ to $h_-$; and $I_4$ is the remaining segment back to $j-r_0$.
Opposite pairs have vertex-set distances
\[
 d(I_1,I_3)\ge R-r_0\ge R/3,\qquad
 d(I_2,I_4)=2r_0\ge R/2-2\ge R/3.
\]
The equality uses that $I_2,I_4$ are subsegments of one geodesic.
Set $q_0=R/3$ and define $1$-Lipschitz vertex functions
\[
 f(v')=\min\{q_0,d(v',I_1)\},\qquad
 g(v')=\min\{q_0,d(v',I_2)\}.
\]
On an oriented edge $(v',v'')$ define the antisymmetric form
\[
 \omega(v',v'')=
 \frac{f(v')+f(v'')}{2}\bigl(g(v'')-g(v')\bigr).
\]
The successive arcs have $f=0$, $g=0$, $f=q_0$, and $g=q_0$,
respectively.  Consequently the integral around the loop is exactly
$q_0^2$: only $I_3$ contributes, with its $g$ values running from zero
to $q_0$.  On any loop of length at most four, subtract the initial value
of $f$ without changing the integral, since the increments of $g$ sum to
zero.  The Lipschitz bounds then bound its absolute integral by $16$.
The integral depends only on the oriented chain, so the assumed short
chain filling gives
\[
 (R/3)^2\le16\cdot12\cdot200R.
\]
Thus $R\le345600<10^6$.  No local finiteness was used: all triangles,
loops and chain decompositions here are finite.  Geodesics between vertices
exist because nonempty sets of integer path lengths have minima.
\end{proof}

\subsection{Avoidance paths and simultaneous overlap}
\label{geo:overlap}

\begin{lemma}[A short path avoiding one center]
\label{geo:avoidance}
Fix a positive integer $L$.  For sufficiently large $\ell$, uniformly in
$Q$ and $i$, if a path in $\widehat X$ of length at most $L$ joins
$u,x\in Q$, with $u$ in layer $i$, and avoids $c_Q$, then
\begin{equation}
 d_Q(u,x)\le24(L+1)L_i.
 \label{geo:avoidance-bound}
\end{equation}
Here the threshold for $\ell$ can depend on $L$, which will subsequently
be an absolute constant.
\end{lemma}

\begin{proof}
Close the given path by a designated two-spoke side through $c_Q$, inflate,
and minimize.  There are at most $L+1$ sides.  We claim every arc occurrence
on the designated adjustable side is shorter than its local $L_j$.
For an opposite bounded-face incidence this is Lemma~\ref{geo:arcs};
an opposite generator incidence has length one.  If the opposite incidence
is on another adjustable side, that side uses a different center because
the original path avoids $c_Q$.  A long arc would align the two source
paths and then identify the copies, as observed after
Proposition~\ref{geo:embeddings}, which is impossible.

It remains to exclude an arc with both incidences on the designated side.
Its edges are exterior bridges.  In the linear order on the designated
side, the interval between the two traversals of such a bridge is an
excursion through the component cut off by that bridge.  All its exterior
positions lie within that interval; in particular it contains no side
break.  Its endpoint group values agree.  Injectivity of $Q$ implies that
its assigned source endpoints agree, so this excursion can be erased from
the adjustable path.  Discard the cut-off component and bridge as well.
The bounded contours in the retained component are unchanged, and no
exterior side break is lost.  This reduces the number of faces, or keeps
it fixed and decreases the number of edges, contrary to minimality.

The designated path therefore has at most $2b\le12(L+1)$ short arc
occurrences.  For clarity their local length bounds do give a uniform
total bound, without presupposing a diameter estimate.  If a prefix has
length $t$, its next initial layer $j$ has $H_j\le H_i+t$, so the next
arc has length at most $L_i+100t/\log_2d$.  Induction over at most
$12(L+1)$ arcs bounds each encountered $H_j/H_i$ by a quantity tending
to one as $\ell\to\infty$.  In particular each arc has length at most
$2L_i$ for sufficiently large $\ell$.  Their total is at most
$24(L+1)L_i$, proving the claim.
\end{proof}

\begin{proposition}[One ball for all earlier intersections]
\label{geo:balls}
There is a well-order of the indexed copies for which the first three
assertions of Theorem~\ref{geo:main} hold.
\end{proposition}

\begin{proof}
Fix one group vertex $o\in X$.  There are countably many indexed copies:
the components of the countable graph $P$ are countable in number and a
map from a fixed component is specified by one group value.  Enumerate
them, and order the centers first by $d_{\widehat X}(c_Q,o)$ and then by
this enumeration.  This is a well-order, since these distances and the
enumeration values are nonnegative integers.  Choose for each $A=c_Q$ one
geodesic from $A$ to $o$, and let its first vertex after $A$ be $u\in Q$.
The choice of $u$ is made once for this copy.

Suppose $x\in Q$ also belongs to a distinct copy with center $B$ satisfying
$d(B,o)\le d(A,o)$.  The path $A,x,B$ is a geodesic of length two,
since distinct centers have no direct edge.  A geodesic from $B$ to $o$
avoids $A$: otherwise its length is at least $2+d(A,o)>d(B,o)$.
Let $k=d(A,o)$.  If $k<D_0+3$, the path from $x$ through $B$ and $o$
to $u$, using the chosen two geodesics, avoids $A$ and has length at most
$1+k+(k-1)=2k$.  If $k\ge D_0+3$, take the vertex $p$ at position
$t=D_0+3$ on the chosen $A,o$ geodesic.  Every vertex of $[A,B]$ is at
distance greater than $D_0$ from $p$, since its distance from $A$ is at
most two.  By Lemma~\ref{geo:thin}, $p$ has a bridge of length at most
$D_0$ to a vertex $q$ of the chosen $B,o$ geodesic.  This bridge avoids
$A$ because $d(p,A)>D_0$.  Furthermore
\[
 d(B,q)=d(B,o)-d(q,o)\le k-(k-t-D_0)=t+D_0.
\]
The path $x$ to $B$, along $[B,q]$, across this bridge to $p$, and back
along $[p,u]$ therefore avoids $A$ and has length at most $2t+2D_0$.
Both cases are bounded by the absolute constant
\begin{equation}
 L^*=2+2(D_0+3)+2D_0=4D_0+8.
 \label{geo:root-path-length}
\end{equation}

Let $i$ be the layer of the fixed $u$, and define
\[
 R_Q=\left\lceil25(L^*+1)L_i\right\rceil,
 \qquad B(Q)=\{v\in Q:d_Q(u,v)\le R_Q\},
\]
with the induced edges.  Lemma~\ref{geo:avoidance} puts every intersection
with an earlier copy into this one ball.  The root $o$, chosen geodesic,
and ball center $u$ have not been changed as the earlier copy varies.

The short chain and thin-triangle arguments needed only $C_*\ge10^4$,
so $D_0$ and $L^*$ have already been fixed independently of any further
increase in $C_*$.  We now fix, for example,
\begin{equation}
 C_*=10^7(L^*+1),\qquad C_\bullet=2\cdot10^4(L^*+1).
 \label{geo:constants}
\end{equation}
There is consequently no circular choice of constants.  Increase $\ell$
to meet all the fixed-constant comparisons made above and below.

An induced ball of radius $R_Q$ with a cycle has a circuit of length at
most $2R_Q+1$: take a shortest-path spanning tree from $u$ in the ball
and any edge not in that tree.  Such a tree exists by choosing a predecessor
at distance one less for each non-root vertex.  Likewise, if a vertex
outside the ball has two edge incidences into it, those incidences and
the tree path between their endpoints give a circuit of length at most
$2R_Q+2$, including the parallel-edge case.  All starting layers involved
differ from $i$ by at most $R_Q+1$, so have $H_j\ge H_i/2$ for large
$\ell$.  Since
$R_Q\le2500(L^*+1)H_i/\log_2d+1$, either circuit would violate
\eqref{geo:girth} with \eqref{geo:constants}.  Thus the ball is a tree,
and every outside vertex has at most one edge incidence to it.  The graph
of edges meeting the ball is precisely this tree with possible leaves
attached, and is a forest.

Finally, its number of vertices is bounded by
\[
 |V(B(Q))|\le(1+6D)^{R_Q+1}.
\]
For large $\ell$, $\log_2(1+6D)\le3\log_2d$, so the logarithm of this
bound is at most
$7500(L^*+1)H_i+6\log_2d\le8000(L^*+1)H_i$.
For every layer $j$ represented in the ball,
$H_j\ge H_i/2$.  Our choice of $C_\bullet$ therefore gives
$|V(B(Q))|\le2^{C_\bullet H_j}$, as required at every vertex in the ball.
\end{proof}

\subsection{Cycle coordinates, the residual forest, and copy incidences}
\label{geo:cycles}

Because all earlier intersections with a copy lie in one tree, a nonzero
cycle in the latest participating copy cannot be canceled by earlier copy
cycles.  The resulting cycle decomposition will control the unused edges
and later supply the
acyclic complex used to exclude torsion.

\begin{proposition}[Direct cycle coordinates]
\label{geo:cycle-directness}
For every field $k$, inclusion of the indexed copies induces an isomorphism
of vector spaces
\begin{equation}
 \bigoplus_Q Z_1(Q;k)\ \longrightarrow\ Z_1(X;k),
 \qquad (c_Q)_Q\longmapsto\sum_Q c_Q.
 \label{geo:cycle-decomposition}
\end{equation}
The direct sum is algebraic, so every tuple used has finite support.
\end{proposition}

\begin{proof}
A finite-support cycle in any graph is a linear combination of finitely
many circuit chains: choose a spanning forest of its finite support graph,
and subtract for each edge outside that forest the corresponding circuit
with its given coefficient.  The remainder is a cycle supported on a finite
forest, and hence is zero by removing leaves.  This proof works over every
field and with loop and parallel edges.

A closed walk in $X$ reads a word trivial in $\Gbase$.  By the presentation
and the van Kampen Lemma it bounds a finite diagram whose faces are assigned
closed source walks in $P$.  The propagated group values give each such
source component its label map into $X$, and therefore its indexed copy.
As oriented edge chains, the exterior is the sum of the bounded face
contours.  Reduction modulo $\operatorname{char} k$ or multiplication by
the circuit coefficients proves that every element of $Z_1(X;k)$ is a
finite sum of copy cycles.  This proves surjectivity.

For injectivity take a finite relation $\sum_Qc_Q=0$ and the latest copy
$Q$ with $c_Q\ne0$.  If an edge has nonzero coefficient in $c_Q$, it has
to occur in an earlier coordinate as well.  Its endpoints then belong
to an earlier copy, and hence both belong to $B(Q)$.  Thus $c_Q$ is a
cycle supported on the tree $B(Q)$, which forces $c_Q=0$.  This
contradiction proves directness.  No claim about independence of infinite
sums has been used.
\end{proof}

\begin{proposition}[Assigned domains and the remaining forest]
\label{geo:forest}
Put
\[
 D_Q=V(Q)\setminus\bigcup_{R<Q}V(R),\qquad
 E_{\mathrm{int}}=\bigcup_Q
 \{e\in E(Q):\text{ both endpoints of }e\text{ lie in }D_Q\}.
\]
The $D_Q$ are pairwise disjoint,
$V(Q)\setminus D_Q\subseteq V(B(Q))$, and the formal graph
$(V(X),E(X)\setminus E_{\mathrm{int}})$ is a forest.  Each edge in
$E_{\mathrm{int}}$ belongs to the displayed set for exactly one $Q$.
\end{proposition}

\begin{proof}
If $R<Q$, a vertex of $R$ cannot belong to $D_Q$; this proves disjointness.
The containment follows from Proposition~\ref{geo:balls}.  Disjointness
also ensures that an edge cannot be internal to two assigned domains.

If the remaining graph contained a circuit, its edge chain over
$\mathbb F_2$ would be a nonzero cycle $c$.  Use
Proposition~\ref{geo:cycle-directness} to write $c=\sum_Qc_Q$ with
finitely many nonzero coordinates, and take their latest copy $Q$.
An edge of $c_Q$ whose endpoints are both outside $B(Q)$ cannot occur
in an earlier coordinate, since that would put its endpoints into the
earlier overlap.  If its coefficient were nonzero it would therefore
have nonzero coefficient in $c$.  But both its endpoints lie in $D_Q$,
so it was deleted, contrary to the support of $c$.  Consequently $c_Q$
is supported on edges with at least one endpoint in $B(Q)$.  These edges
form a forest by Proposition~\ref{geo:balls}, forcing $c_Q=0$, a
contradiction.
\end{proof}

\begin{proposition}[The exact incidence count]
\label{geo:incidences}
Let
\[
 \mathcal I=\{(Q,x):Q\text{ an indexed copy},\ x\in V(Q)\}.
\]
There is a canonical bijection
\begin{equation}
 \mathcal I\longrightarrow\Gbase\times V(Y),\qquad
 (Q,x)\longmapsto(x,\operatorname{type}_Q(x)),
 \label{geo:incidence-bijection}
\end{equation}
where the type is the base vertex under the unique source vertex of the
embedding.  It intertwines left translation of $\Gbase$ with left
translation of the first coordinate, with the base type fixed.
In particular, for a fixed $x\in\Gbase$ and $v\in V(Y)$ there is exactly
one incidence at $x$ of type $v$, even when $P$ is disconnected.
\end{proposition}

\begin{proof}
Embedding makes the source vertex of $x$ unique within a representative
$(C,\phi)$; deck alignment preserves its base vertex, so the type is
well-defined.  For given $(x,v)$ choose the lift $(v,1_F)$ and its component
$C$, and choose the unique label map from $C$ sending that lift to $x$.
This constructs an incidence.  Conversely a representative of any such
incidence has its source vertex equal to $(v,f)$ for some $f\in F$.
Apply $\delta_{f^{-1}}$ and change the map by precomposition with
$\delta_f$.  This is the equivalence defining indexed copies and brings
that source vertex to $(v,1_F)$.  Its component is now $C$, and its map
is the unique one just constructed.  This proves uniqueness.  Left
translation changes $x$ and not the source data, giving equivariance.
\end{proof}

The next two estimates control different costs of removing overlaps.  The
first sums over one trimming ball, allowing repeated lifts of a base
vertex.  The second sums the losses from all copies at one fixed Cayley
vertex, retaining the size of each trimming ball.  If nonnegative numbers
$a_v$ are assigned to base vertices and
$a_Q(x)=a_{\operatorname{type}_Q(x)}$, put
\[
 \mathcal S_a=\sum_{i\in\mathbb Z}\mathfrak w_i
                   \sum_{v\text{ in layer }i}a_v.
\]
Then, whenever this sum is finite,
\begin{align}
 \sum_{x\in B(Q)}a_Q(x)&\le\mathcal S_a,
 \label{geo:ball-weighted-sum}\\
 \sum_{Q:\,x\in V(Q)\setminus D_Q}|V(B(Q))|a_Q(x)
          &\le\mathcal S_a\qquad(x\in\Gbase).
 \label{geo:lost-weighted-sum}
\end{align}
For the first inequality, a base vertex in layer $i$ occurs among the
lifts in $B(Q)$ at most $|V(B(Q))|\le\mathfrak w_i$ times.  For the
second, every incidence being summed has $x\in B(Q)$, so its ball size
is at most the weight of its type's layer; the types occur at most once
by Proposition~\ref{geo:incidences}.  In particular the equivariant
incidence Hilbert space is canonically
\[
 \bigoplus_Q\ell^2(V(Q))\otimes\mathbb C^N
 \cong \ell^2(\Gbase)\otimes\ell^2(V(Y))\otimes\mathbb C^N.
\]
For a deck-equivariant family of transported projections with base
diagonal matrix sums $a_v$, its direct sum has unnormalized identity
diagonal sum $\sum_v a_v$.  This is a statement about all indexed copies;
the well-order and assigned domains need not be equivariant.

\subsection{Cone homology and exclusion of torsion}
\label{geo:torsion-subsection}

\begin{lemma}[Stabilizers of indexed copies]
\label{geo:stabilizers}
The stabilizer in $\Gbase$ of an indexed copy embeds into the subgroup
of $F$ preserving any representative source component.  In particular
that stabilizer has no nonidentity finite-order element.
\end{lemma}

\begin{proof}
Fix a representative $(C,\phi)$.  If $h$ stabilizes the indexed copy, its
definition means that there is a deck transformation $\delta_a$ preserving
$C$ such that
\[
 h\phi=\phi\delta_a\quad\text{on }C.
\]
It is unique because $\phi$ is injective and the deck action is free on
vertices.  If $h_1,h_2$ correspond to $a_1,a_2$, the displayed equality
gives $(h_1h_2)\phi=\phi\delta_{a_1a_2}$, so the correspondence is a
homomorphism.  If $a=1_F$ it implies that $h$ fixes a group vertex, so
$h=1_{\Gbase}$.  Hence it is injective.  Torsion-freeness follows from
that of $F$.  This argument does not assert that the full stabilizer is
trivial and does not require a component of $P$ to be the entire cover.
\end{proof}

Let $Z$ be the two-dimensional cellular complex obtained by attaching the
graph cone on each indexed copy to $X$.  Its one-skeleton is $\widehat X$.
For each positive edge $e:s(e)\to t(e)$ in $Q$, attach one triangular
cell $\tau_{Q,e}$ with boundary
\begin{equation}
 \partial\tau_{Q,e}=e+\sigma_{Q,s(e)}-\sigma_{Q,t(e)},
 \label{geo:cone-boundary}
\end{equation}
where $\sigma_{Q,x}$ is the oriented spoke from $c_Q$ to $x$.  Edges
belonging to different indexed copies give distinct triangular cells,
even when their base edge in $X$ is the same.  This is a cellular
attachment, so repeated incidences would be allowed; the embeddings
already proved in fact exclude loops in individual copies in this
construction.

\begin{lemma}[Acyclicity of the cone complex]
\label{geo:acyclic}
For every field $k$, the augmented cellular chain complex
\begin{equation}
 0\longrightarrow C_2(Z;k)\xrightarrow{\partial_2}C_1(Z;k)
 \xrightarrow{\partial_1}C_0(Z;k)\xrightarrow{\varepsilon}k
 \longrightarrow0
 \label{geo:acyclic-complex}
\end{equation}
is exact.
\end{lemma}

\begin{proof}
The graph $X$ is connected and each cone is attached to it, so exactness
at $C_0$ and $k$ is immediate.  To prove exactness at $C_1$, start with
a finite chain cycle in $\widehat X$.  For each center its spoke
coefficients sum to zero, since its boundary coefficient is zero.  Only
finitely many spoke coefficients are nonzero.  As its copy is connected,
a finite zero-sum chain of these copy vertices is the boundary of a
finite edge chain in that copy: join each supporting vertex to a fixed
one by a finite path and take the corresponding linear combination.
Formula~\eqref{geo:cone-boundary} therefore allows addition of a finite
combination of cone triangle boundaries to eliminate all spoke terms.
The resulting cycle lies in $X$.  By
Proposition~\ref{geo:cycle-directness}, it is a finite sum of copy cycles,
and each copy cycle is the boundary of the same linear combination of
its cone triangles, since its spoke terms cancel.  Thus every one-cycle
is a boundary.

For injectivity at $C_2$, a finite sum of triangles determines, for each
copy, the finite edge chain $c_Q$ of its triangle coefficients.  Vanishing
of the boundary on that copy's spokes says exactly $\partial c_Q=0$,
by \eqref{geo:cone-boundary}.  The remaining boundary in $X$ is
$\sum_Qc_Q$.  Its vanishing, together with directness in
\eqref{geo:cycle-decomposition}, implies $c_Q=0$ for every $Q$.
Each triangle corresponds to a distinct oriented edge in that copy,
so all triangle coefficients vanish.  There are no cells in higher
degrees.  All these arguments concern finite-support chains, also for
infinite copies and an infinite family of cones.
\end{proof}

\begin{proposition}[Torsion-freeness]
\label{geo:torsion}
The finitely generated group $\Gbase$ in \eqref{geo:presentation} is
torsion-free.
\end{proposition}

\begin{proof}
Left translation acts cellularly on $Z$.  A nontrivial finite cyclic
subgroup acts freely on its cells.  Indeed stabilizing an original group
vertex is impossible; stabilizing a formal generator edge preserves its
positive orientation and fixes its initial vertex.  Stabilizing a center,
a spoke, or a triangle entails stabilizing the indexed copy given by
that center.  Lemma~\ref{geo:stabilizers} excludes any nonidentity
finite-order element doing so.  This verifies cell freeness, not merely
freeness at generic points, and avoids an assumption about edge inversions.

If $\Gbase$ had a nonidentity torsion element, one of its powers would
generate a subgroup $C_p$ of prime order $p$.  Over $k=\mathbb F_p$,
choose orientations and orbit representatives of its cells.  Because
the action on cells is free, every $C_j(Z;k)$ is a free module over
$A=k[C_p]$, possibly of infinite rank.  By Lemma~\ref{geo:acyclic},
\eqref{geo:acyclic-complex} would be a free resolution of the trivial
$A$-module $k$ of length two.

The obstruction is the standard periodic-resolution argument for finite
cyclic groups; see Brown \cite[Examples~1.3 and~2.4(e)]{Brown2010}.
We give it here over $k$.  Writing $t$
for a generator of $C_p$ and $z=t-1$, one has
$A=k[z]/(z^p)$, and $k=A/(z)$.  The complex
\begin{equation}
 \cdots\xrightarrow{z}A\xrightarrow{z^{p-1}}A
 \xrightarrow{z}A\xrightarrow{z^{p-1}}A
 \xrightarrow{z}A\longrightarrow k\longrightarrow0
 \label{geo:periodic-resolution}
\end{equation}
is exact: in the basis $1,z,\ldots,z^{p-1}$, the kernel of multiplication
by $z$ is $(z^{p-1})$, and the kernel of multiplication by $z^{p-1}$ is
$(z)$.  Tensoring it over $A$ with $k$ makes every differential zero,
and leaves one copy of $k$ in every nonnegative degree.

Any two free resolutions of $k$ give the same homology after this tensor
operation.  Here is the comparison argument, including infinite free
ranks.  Lift the identity of $k$ to a chain map between the resolutions:
on each free basis vector in degree zero choose a preimage of its
augmentation; inductively, the image already assigned to a boundary
is a cycle, and exactness in the other resolution supplies its preimage.
Extend linearly on the free basis.  Construct a map in the opposite
direction in the same way.  Their composites lift the identity, and
their differences from the identity are null-homotopic: first lift the
degree-zero difference, whose augmentation vanishes, and then lift the
difference in each next degree after subtracting the already chosen
homotopy term.  At each step it is a cycle by the chain-map identity.
These chain homotopies persist on tensoring with $k$, proving the claimed
equality of homology.  Applied to the length-two resolution and
\eqref{geo:periodic-resolution}, it would make its homology both zero
and equal to $k$ in degrees greater than two.  This contradiction rules
out $C_p$, and hence all nonidentity torsion.

Finite generation follows directly from the finiteness of $S$ in
\eqref{geo:presentation}.  Combining the preceding results proves
Theorem~\ref{geo:main} in full.
\end{proof}
